% =====================================================================
%  courses/aritmetica/semana_01_logica/teoria.tex — teoría: conectivos lógicos y tabla de verdad
% ---------------------------------------------------------------------
%  Curso: Aritmética / Lógica | Semana: 01
%  MÓDULO: sin \documentclass ni \begin{document}
% =====================================================================
	
	% --- PORTADA DE CLASE (usa \portadaclase definido en preambulo_comun) ------
	\portadaclase%
	{Conectivos Lógicos y Tabla de Verdad}%  #1 Tema
	{01}%                                    #2 Semana
	{I}%                                     #3 N° Unidad
	{Fundamentos Numéricos}%                 #4 Nombre Unidad
	
	% --- MOTIVACIÓN ------------------------------------------------------------
	\section*{¿Por qué estudiar este tema?}
	\begin{tcolorbox}[
		enhanced, colback=AzulClaro!30, colframe=AzulMedio,
		arc=3pt, boxrule=0.6pt, left=10pt, right=10pt
		]
		\textit{Los conectivos lógicos y las tablas de verdad son base obligatoria en los
			exámenes de admisión de la UNI, UNMSM, UNSCH, PUCP y Villarreal. Dominar el
			método \textbf{backdoor} y la clasificación de matrices (tautología, contradicción,
			contingencia) permite resolver este bloque completo en tiempo mínimo.}
	\end{tcolorbox}
	\vspace{0.4cm}
	
	% --- SECCIÓN 1: LÓGICA PROPOSICIONAL ---------------------------------------
	\section{Lógica Proposicional}
	
	\subsection{Definiciones fundamentales}
	
	\begin{cajadef}
		\textbf{Lógica proposicional:} Rama de la lógica que estudia las
		\textbf{proposiciones}, las relaciones entre ellas y la función de las
		variables proposicionales y los conectivos lógicos.
	\end{cajadef}
	
	\begin{cajadef}[2]
		\textbf{Proposición lógica} (enunciado cerrado): Todo enunciado declarativo
		con un único valor de verdad: \textbf{Verdadero~(V)} o \textbf{Falso~(F)}.
		Se representa con letras minúsculas: $p$, $q$, $r$, $s$\ldots
	\end{cajadef}
	
	\subsection{Clasificación de proposiciones}
	
	\begin{cajaprop}[Tipos]
		\begin{itemize}
			\item \textbf{Simples (atómicas):} Sin conectivos lógicos.
			Ejemplo: \textit{``Carlos es travieso''}.
			\item \textbf{Compuestas (moleculares):} Dos o más simples unidas por un
			conectivo. Ejemplo: \textit{``Carlos es travieso \textbf{y} despistado''}.
		\end{itemize}
	\end{cajaprop}
	
	\begin{cajaadvert}
		\textbf{Error frecuente:} Confundir sujeto compuesto con proposición compuesta.\\
		``Carlos y Richard van al cine'' $\to$ sujeto compuesto, acción única $\to$
		proposición \textbf{simple}.\\[4pt]
		\textbf{Regla:} Una proposición es compuesta si se puede separar en dos oraciones
		con valor de verdad propio e independiente.
	\end{cajaadvert}
	
	% --- SECCIÓN 2: CONECTIVOS LÓGICOS -----------------------------------------
	\section{Conectivos Lógicos}
	
	\subsection{Conjunción }
	
	\begin{cajadef}
		\textbf{Conjunción:} Conectiva \textbf{y} (también: \textit{pero, e, sin
			embargo}). Operador: $\wedge$. Proposición resultante: $p \wedge q$.
	\end{cajadef}
	
	\begin{cajaprop}[Conjunción]
		$p \wedge q$ es \textbf{V} únicamente cuando \textbf{ambas} proposiciones son V.
		\[
		\begin{array}{|c|c||c|}
			\hline
			\rowcolor{AzulClaro!60}
			p & q & p \wedge q \\
			\hline
			V & V & \mathbf{V} \\
			V & F & F \\
			F & V & F \\
			F & F & F \\
			\hline
		\end{array}
		\]
	\end{cajaprop}
	
	\subsection{Disyunción débil ($\vee$)}
	
	\begin{cajadef}
		\textbf{Disyunción débil:} Conectiva \textbf{o} (también: \textit{u}).
		Operador: $\vee$. Proposición: $p \vee q$.
	\end{cajadef}
	
	\begin{cajaprop}[Disyunción débil]
		$p \vee q$ es \textbf{F} únicamente cuando \textbf{ambas} son F.
		\[
		\begin{array}{|c|c||c|}
			\hline
			\rowcolor{AzulClaro!60}
			p & q & p \vee q \\
			\hline
			V & V & V \\
			V & F & V \\
			F & V & V \\
			F & F & \mathbf{F} \\
			\hline
		\end{array}
		\]
	\end{cajaprop}
	
	\subsection{Disyunción fuerte ($\Delta$)}
	
	\begin{cajadef}
		\textbf{Disyunción fuerte (exclusiva):} Conectiva \textbf{o\ldots o / o
			bien\ldots o bien}. Operador: $\Delta$. Proposición: $p \,\Delta\, q$.
	\end{cajadef}
	
	\begin{cajaprop}[Disyunción fuerte]
		$p \,\Delta\, q$ es \textbf{V} cuando exactamente \textbf{una} proposición
		es V (valores distintos entre sí).
		\[
		\begin{array}{|c|c||c|}
			\hline
			\rowcolor{AzulClaro!60}
			p & q & p \,\Delta\, q \\
			\hline
			V & V & F \\
			V & F & \mathbf{V} \\
			F & V & \mathbf{V} \\
			F & F & F \\
			\hline
		\end{array}
		\]
	\end{cajaprop}
	
	\subsection{Condicional ($\to$)}
	
	\begin{cajadef}
		\textbf{Condicional:} Conectiva \textbf{si\ldots entonces} (también:
		\textit{implica, por lo tanto}). Operador: $\to$. Proposición: $p \to q$.
		$p$ = antecedente; $q$ = consecuente.
	\end{cajadef}
	
	\begin{cajaprop}[Condicional]
		$p \to q$ es \textbf{F} únicamente cuando $p = V$ y $q = F$.
		\[
		\begin{array}{|c|c||c|}
			\hline
			\rowcolor{AzulClaro!60}
			p & q & p \to q \\
			\hline
			V & V & V \\
			V & F & \mathbf{F} \\
			F & V & V \\
			F & F & V \\
			\hline
		\end{array}
		\]
	\end{cajaprop}
	
	\begin{cajatip}
		Mnemotecnia: el condicional es una \textbf{promesa}. Solo se rompe (F) cuando
		se prometió ($p=V$) y no se cumplió ($q=F$). Sin promesa ($p=F$) no hay
		incumplimiento $\Rightarrow$ siempre V.
	\end{cajatip}
	
	\subsection{Bicondicional ($\leftrightarrow$)}
	
	\begin{cajadef}
		\textbf{Bicondicional:} Conectiva \textbf{si y solo si} (también:
		\textit{entonces y solo entonces}). Operador: $\leftrightarrow$ (también $\equiv$).
		Proposición: $p \leftrightarrow q$.
	\end{cajadef}
	
	\begin{cajaprop}[Bicondicional]
		$p \leftrightarrow q$ es \textbf{V} cuando ambas tienen el
		\textbf{mismo valor de verdad}.
		\[
		\begin{array}{|c|c||c|}
			\hline
			\rowcolor{AzulClaro!60}
			p & q & p \leftrightarrow q \\
			\hline
			V & V & V \\
			V & F & F \\
			F & V & F \\
			F & F & V \\
			\hline
		\end{array}
		\]
	\end{cajaprop}
	
	\subsection{Negación ($\sim$)}
	
	\begin{cajadef}
		\textbf{Negación:} Conectiva \textbf{no} (también: \textit{ni, es falso que}).
		Operador: $\sim$ (también $\neg$). Proposición: $\sim p$.
	\end{cajadef}
	
	\begin{cajaprop}[Negación]
		$\sim p$ invierte siempre el valor de $p$.
		\[
		\begin{array}{|c||c|}
			\hline
			\rowcolor{AzulClaro!60}
			p & \sim p \\
			\hline
			V & F \\
			F & V \\
			\hline
		\end{array}
		\]
	\end{cajaprop}
	
	% --- SECCIÓN 3: TABLAS DE VERDAD -------------------------------------------
	\section{Tablas de Valores de Verdad}
	
	\subsection{Definición y número de filas}
	
	\begin{cajadef}
		\textbf{Tabla de verdad:} Matriz que lista todas las combinaciones posibles
		de valores de las variables y el valor resultante de la proposición compuesta.
		La columna del operador principal se llama \textbf{matriz principal}.
	\end{cajadef}
	
	\begin{cajaprop}[Número de filas]
		Con $n$ variables proposicionales:
		\[
		\text{N}^{\circ}\text{ de filas} = 2^{n}
		\]
		\vspace{-6pt}
		\begin{itemize}
			\item 1 variable: $2^1 = 2$ filas
			\item 2 variables: $2^2 = 4$ filas
			\item 3 variables: $2^3 = 8$ filas
		\end{itemize}
	\end{cajaprop}
	
	\subsection{Jerarquía de conectivos}
	
	\begin{cajaformula}
		\begin{center}
			\renewcommand{\arraystretch}{1.4}
			\begin{tabular}{@{}clc@{}}
				\toprule
				\rowcolor{AzulClaro!60}
				\textbf{Prioridad} & \textbf{Conectivo} & \textbf{Símbolo} \\
				\midrule
				1\textsuperscript{o} (mayor) & Negación         & $\sim$            \\
				\rowcolor{GrisClaro}
				2\textsuperscript{o}          & Conjunción       & $\wedge$          \\
				3\textsuperscript{o}          & Disyunción débil & $\vee$            \\
				\rowcolor{GrisClaro}
				4\textsuperscript{o}          & Disyunción fuerte& $\Delta$          \\
				5\textsuperscript{o}          & Condicional      & $\to$             \\
				\rowcolor{GrisClaro}
				6\textsuperscript{o} (menor)  & Bicondicional    & $\leftrightarrow$ \\
				\bottomrule
			\end{tabular}
		\end{center}
	\end{cajaformula}
	
	\begin{cajaadvert}
		\textbf{Error:} Evaluar de izquierda a derecha sin respetar jerarquía.\\
		\textbf{Corrección:} Identificar el operador principal (menor prioridad fuera
		de paréntesis) y evaluarlo \textbf{al final}.
	\end{cajaadvert}
	
	\subsection{Ejemplo completo de tabla}
	
	Proposición: $(p \wedge q) \to (p \,\Delta\, q)$ \quad
	(Orden: \textcircled{1} $p \wedge q$, \textcircled{2} $p\,\Delta\,q$,
	\textcircled{3} $\to$)
	
	\[
	\begin{array}{|cc|c|c||c|}
		\hline
		\rowcolor{AzulClaro!60}
		p & q & \textcircled{1}\;p \wedge q & \textcircled{2}\;p\,\Delta\,q
		& \textcircled{3}\;(p \wedge q) \to (p\,\Delta\,q) \\
		\hline
		V & V & V & F & F \\
		V & F & F & V & V \\
		F & V & F & V & V \\
		F & F & F & F & V \\
		\hline
	\end{array}
	\]
	
	\noindent Matriz principal: \textbf{FVVV} $\Rightarrow$ proposición
	\textbf{contingente}.
	
	% --- SECCIÓN 4: CLASES DE MATRICES -----------------------------------------
	\section{Clases de Matrices Principales}
	
	\begin{cajaprop}[Tautología]
		Todos los valores de la matriz principal son \textbf{V}.\\
		Notación: $\models \varphi$ \quad (verdad lógica).
	\end{cajaprop}
	
	\begin{cajaprop}[Contradicción]
		Todos los valores de la matriz principal son \textbf{F}.\\
		Notación: $\varphi \equiv \bot$ \quad (también: \textit{antología}).
	\end{cajaprop}
	
	\begin{cajaprop}[Contingencia]
		La matriz principal contiene \textbf{al menos un V y un F}.
	\end{cajaprop}
	
	% --- SECCIÓN 5: EJEMPLOS RESUELTOS -----------------------------------------
	\section{Ejemplos resueltos}
	
	\begin{cajaejemplo}[1 — Identificar proposición compuesta]
		\textbf{Enunciado:} ¿Cuál es una proposición compuesta?\\[4pt]
		a)~``Agripino y Cesarina son hermanos.'' \quad
		e)~``Daniel es profesor y Rosa es escritora.''
		
		\medskip
		\noindent\textbf{Método:} Análisis estructural.
		
		\medskip
		\noindent (a): sujeto compuesto, una sola acción $\to$ \textbf{simple}.\\
		(e): $p$: ``Daniel es profesor'' y $q$: ``Rosa es escritora'' $\to$
		dos proposiciones independientes $\to$ \textbf{compuesta}: $p \wedge q$.
		
		\begin{tcolorbox}[colback=AzulClaro!50, colframe=AzulPizarra,
			arc=2pt, boxrule=0.7pt]
			\[\boxed{(e)\;p \wedge q:\;\text{``Daniel es profesor''}
				\wedge \text{``Rosa es escritora''}}\]
		\end{tcolorbox}
	\end{cajaejemplo}
	
	\begin{cajaejemplo}[2 — Simbolización]
		\textbf{Enunciado:} Simboliza: \textit{``Si Daniel y Agripina juegan fútbol,
			entonces Margarito será el árbitro''}.
		
		\medskip
		$p$: Daniel juega fútbol,\quad $q$: Agripina juega fútbol,\quad
		$r$: Margarito será árbitro.
		
		\begin{align*}
			\text{Antecedente:} &\quad p \wedge q \\
			\text{Consecuente:} &\quad r \\
			\text{Resultado:}   &\quad (p \wedge q) \to r
		\end{align*}
		
		\begin{tcolorbox}[colback=AzulClaro!50, colframe=AzulPizarra,
			arc=2pt, boxrule=0.7pt]
			\[\boxed{(p \wedge q) \to r}\]
		\end{tcolorbox}
	\end{cajaejemplo}
	
	\begin{cajaejemplo}[3 — Método backdoor]
		\textbf{Enunciado:} La proposición $(\sim p \wedge q) \to [(p \vee r) \vee t]$
		es \textbf{falsa}. Halla $p$, $q$, $r$, $t$.
		
		\medskip
		\noindent Un condicional es F solo cuando ant.~$= V$ y cons.~$= F$:
		
		\begin{align*}
			\text{Consecuente F:}&\quad
			(p \vee r) \vee t = F \Rightarrow p=F,\;r=F,\;t=F \\
			\text{Antecedente V:}&\quad
			\sim p \wedge q = V \Rightarrow \underbrace{\sim p=V}_{\Rightarrow p=F\;\checkmark}
			\text{ y } q=V
		\end{align*}
		
		\begin{tcolorbox}[colback=AzulClaro!50, colframe=AzulPizarra,
			arc=2pt, boxrule=0.7pt]
			\[\boxed{p \equiv F,\quad q \equiv V,\quad r \equiv F,\quad t \equiv F}\]
		\end{tcolorbox}
		
		\begin{cajatip}
			\textbf{Backdoor:} trabaja siempre desde el conectivo principal hacia adentro.
			Condicional falso $\Rightarrow$ ant.~V y cons.~F.
			Disyunción falsa $\Rightarrow$ ambos miembros F. Estas dos reglas resuelven
			el 90\,\% de los ejercicios de lógica en admisión.
		\end{cajatip}
	\end{cajaejemplo}
	
	% --- RESUMEN FINAL ---------------------------------------------------------
	\section{Resumen del tema}
	
	\begin{cajaformula}
		\begin{center}
			\renewcommand{\arraystretch}{1.5}
			\begin{tabular}{@{}p{3.8cm}cp{6.8cm}@{}}
				\toprule
				\rowcolor{AzulClaro!60}
				\textbf{Conectivo} & \textbf{Símbolo} & \textbf{¿Cuándo es FALSO?} \\
				\midrule
				Conjunción & $p \wedge q$ & Cuando al menos uno es F \\
				\rowcolor{GrisClaro}
				Disyunción débil & $p \vee q$ & Solo cuando ambos son F \\
				Disyunción fuerte & $p\,\Delta\,q$ & Cuando ambos son iguales (VV o FF) \\
				\rowcolor{GrisClaro}
				Condicional & $p \to q$ & Solo cuando $p=V$ y $q=F$ \\
				Bicondicional & $p \leftrightarrow q$ & Cuando tienen valores distintos \\
				\rowcolor{GrisClaro}
				Negación & $\sim p$ & Cuando $p$ es V \\
				\bottomrule
			\end{tabular}
		\end{center}
	\end{cajaformula}
	
	\begin{cajatip}
		Memoriza la \textbf{excepción de cada conectivo}. El condicional con
		antecedente falso (siempre V) y el bicondicional con ambos falsos (también V)
		son los casos que más sorprenden en el examen.
	\end{cajatip}
	
	% --- FIN — teoria.tex ------------------------------------------------------